\bookchapter{Kapitel 52: Die vier Grundkomponenten und der Activity-Lebenszyklus}

\chapterintro{Wenn Compose das “Gesicht” der App ist, dann ist dieses Kapitel das “Skelett”. Android kennt vier Grundkomponenten, aus denen jede App zusammengesetzt ist; und die wichtigste davon, die Activity, hat eine “Lebenszyklus-Karte”, die jeder Android-Entwickler ins Gehirn gebrannt haben muss. Nach diesem Kapitel verstehst du so alte Rätsel wie “Warum ist mein App-Zustand weg, wenn ich kurz in den Hintergrund wechsle und zurückkomme?”}

\booksection{Die vier Grundkomponenten: Androids tragende Säulen}

Android definiert offiziell vier Grundkomponenten, die jeweils eine Art von Arbeit übernehmen:

\begin{center}
\begin{tabularx}{\linewidth}{|X|X|X|}
\hline
\textbf{Komponente} & \textbf{Aufgabe} & \textbf{Zum Vergleich} \tabularnewline
\hline
Activity & Eine sichtbare Oberfläche & Die Schaufensterfront des Ladens – sichtbar und bedienbar \tabularnewline
Service & Der "Unsichtbare", der im Hintergrund schuftet & Der Lagerarbeiter – Musik abspielen, Dateien herunterladen, ganz ohne Oberfläche \tabularnewline
BroadcastReceiver & System-Broadcasts empfangen & Die Pförtnerloge – der Lautsprecher für "Akku ist leer" und "Netz ist weg" \tabularnewline
ContentProvider & Daten app-übergreifend teilen & Der öffentliche Aktenschrank – z. B. Kontakte für andere Apps freigeben \tabularnewline
\hline
\end{tabularx}
\end{center}

\begin{center}\small\color{gray}Die vier Grundkomponenten von Android auf einen Blick\end{center}

Als Einsteiger verbringst du 95 \% deiner Zeit mit der Activity; Service und BroadcastReceiver brauchst du für Musikwiedergabe und Push-Nachrichten; ContentProvider nutzt du am seltensten (das System kapselt die gebräuchlichen Daten bereits). Und beachte: Alle vier Komponenten müssen sich im AndroidManifest.xml “anmelden” – ohne Registrierung findet das System sie nicht.

\booksection{Der Activity-Lebenszyklus: die erste “Schatzkarte” der Android-Entwicklung}

Eine App auf dem Telefon kann jederzeit unterbrochen werden: Ein Anruf kommt, der Nutzer wechselt kurz zu WeChat, das System beendet dich wegen Speichermangels … Das Leben einer Activity ist eine Rundreise auf dieser Karte:

onCreate (Geburt: Initialisierung)

onStart (sichtbar)

onResume (im Vordergrund aktiv)

onPause (teilweise verdeckt)

onStop (ganz unsichtbar)

onDestroy (zerstört)

onRestart / onStart / onResume: beim Zurückkehren in den Vordergrund umgekehrt

\begin{center}\small\color{gray}Abbildung 4: Der Activity-Lebenszyklus (vorwärts: obere Reihe Geburt untere Reihe Abschied; auf dem Rückweg genau umgekehrt)\end{center}

\begin{center}
\begin{tabularx}{\linewidth}{|X|X|X|}
\hline
\textbf{Callback} & \textbf{Wann wird er ausgelöst?} & \textbf{Typischer Einsatz} \tabularnewline
\hline
onCreate & Beim Erzeugen der Activity (einmal) & Oberfläche laden, Daten initialisieren \tabularnewline
onStart & Die Oberfläche wird sichtbar & Ein wenig leichtgewichtige Logik starten \tabularnewline
onResume & Hat den Fokus, ist bedienbar & Animationen starten, Wiedergabe fortsetzen \tabularnewline
onPause & Teilweise verdeckt oder im Begriff, verlassen zu werden & Animationen pausieren, Entwurf speichern \tabularnewline
onStop & Vollständig unsichtbar & Schwere Ressourcen freigeben, Downloads pausieren \tabularnewline
onDestroy & Unmittelbar vor dem Entfernen & Die letzten Aufräumarbeiten \tabularnewline
\hline
\end{tabularx}
\end{center}

\begin{center}\small\color{gray}Typische Einsätze der sechs Lifecycle-Callbacks\end{center}

\textbf{Stolperfalle: Eine Drehung des Bildschirms “tötet” die Activity}\par

Der klassischste Stolperstein: Man dreht den Bildschirm, und der ganze Text im Eingabefeld ist weg. Grund: Eine Bildschirmdrehung zerstört die Activity standardmäßig und baut sie neu auf (durchlaufen von onDestroy onCreate). Lösung: Den Zustand mit rememberSaveable sichern (in Compose der Ersatz für remember) – es holt die Werte beim Neuaufbau automatisch zurück. In den Kapiteln 6 und 7 wirst du es immer wieder brauchen.

\booksection{Intent: das “Ticket” für den Verkehr in der App und nach außen}

Der Intent (die “Absicht”) ist das “Fahrticket” für die Kommunikation zwischen Komponenten. Ob du von einer Seite zur anderen springen, den Browser öffnen oder die Kamera aufrufen willst – immer musst du “ein Ticket lösen und einsteigen”:

\begin{lstlisting}[language=Kotlin]
// Variante 1: expliziter Intent – klar angegeben, zu welcher Seite
val intent = Intent(this, DetailActivity::class.java)
intent.putExtra("name", "Max")     // nimmt "Gepäck" mit (die Daten)
startActivity(intent)

// Variante 2: impliziter Intent – du sagst nur, was zu tun ist; das System sucht, wer es tut
val web = Intent(Intent.ACTION_VIEW, Uri.parse("https://developer.android.com"))
startActivity(web)   // das System öffnet den Browser (oder lässt dich wählen)
\end{lstlisting}

Beide Varianten haben ihren Zweck: Explizit für Seitensprünge innerhalb der eigenen App; implizit, um Fähigkeiten des Systems oder anderer Apps zu nutzen (Webseite öffnen, anrufen, teilen). In Compose-Projekten holst du in einer Composable mit LocalContext.current den Context und kannst dann startActivity aufrufen.

\booksection{Mehrere Seiten und Navigation}

Echte Apps haben mehrere Seiten. Die offiziell empfohlene Navigationslösung für Compose ist Navigation Compose (Bibliothek androidx.navigation):

\begin{lstlisting}[language=Kotlin]
val navController = rememberNavController()

NavHost(navController, startDestination = "home") {
    composable("home") { HomeScreen(navController) }
    composable("detail/{name}") { entry ->
        val name = entry.arguments?.getString("name")
        DetailScreen(name ?: "Unbekannt")
    }
}

// um innerhalb der Ansicht zu navigieren:
navController.navigate("detail/Max")     // mit Parameter
navController.popBackStack()              // zurück zur vorherigen Seite
\end{lstlisting}

\begin{tcbinfo}{Tipp: Und wenn es zu viele Seiten werden?}
Sobald das Projekt wächst, schreib alle Seitenrouten gebündelt in eine Datei (z. B. AppNavGraph.kt) und verwalte sie zentral – lass die Logik “Wohin springe ich?” nicht auf alle Seiten verstreuen. Das ist der erste Schritt von “läuft irgendwie” zu “wartbar”.

\end{tcbinfo}

\booksection{Nebenbei erklärt: Was ist ein Fragment (und warum lehrt dieses Buch es nicht)}

In alten Tutorials wirst du unweigerlich auf Fragment stoßen – einen “wiederverwendbaren Oberflächen-Schnipsel” innerhalb einer Activity, in der alten Mehrseiten-Architektur der absolute Star. Im Compose-Zeitalter reichen aber Composable-Funktionen + Navigation-Routen für Seiten, und auch offiziell ist Fragment zurückgestuft worden. Deshalb taucht es in diesem Buch nicht im Hauptweg auf.

Mein Rat: In neuen Projekten Fragment nicht aktiv einsetzen; aber wenn du alten Code liest oder im Vorstellungsgespräch danach gefragt wirst, reicht zu wissen: “Ein Fragment ist ein Oberflächen-Schnipsel, das in einer Activity steckt.” Lernmaterialien immer bevorzugt ab 2023 wählen – die “alten Kalenderblätter” des Android-Stacks veralten schneller, als man denkt.

\begin{tcbpractice}{Selbst ausprobieren: Dem Zähler ein "Gedächtnis" verpassen}
Bau den Zähler aus Kapitel 51 um: (1) Ersetze remember durch rememberSaveable und prüfe nach einer Bildschirmdrehung, ob die Zahl überlebt hat; (2) füge einen “Über”-Button hinzu, der mit Navigation auf eine zweite Seite springt und dort die Versionsnummer zeigt; (3) lass in onCreate bzw. onPause beider Seiten Logausgaben schreiben (Log.d) und beobachte in Logcat die Reihenfolge der Lifecycle-Aufrufe. Nach diesen drei Schritten ist der Lebenszyklus keine Theorie mehr.

\end{tcbpractice}

\begin{tcbsummary}{Zusammenfassung des Kapitels}
\begin{itemize}
  \item Die vier Grundkomponenten: Activity (Schaufenster), Service (Hintergrund), BroadcastReceiver (Lautsprecher), ContentProvider (Aktenschrank)
  \item Sechs Lifecycle-Callbacks: onCreate onStart onResume onPause onStop onDestroy
  \item Eine Bildschirmdrehung zerstört die Activity und baut sie neu: Mit rememberSaveable den Zustand retten
  \item Intent ist das "Ticket": explizit zur eigenen Seite, implizit für System-Fähigkeiten
  \item Mehrere Seiten mit Navigation Compose: NavHost + composable-Routen
\end{itemize}

\end{tcbsummary}
